\section{Uniform original-metric error and the two finite singular values}
The exact remainder MR14 also gives bounds with constants
independent of $w$ on a specified interval. Keep $g>0$ unchanged,
put $G_0=G_{U,x}$ with all its original entries from RS24--25,
and define
\begin{equation}
 \begin{gathered}
 \kappa_{G_0}=
   \sqrt{\operatorname{tr}G_0\operatorname{tr}G_0^{-1}},\\
 B(g)=\frac9{256}+\frac3{256g^2}+\frac1{4g^4}
             +\frac3{16g^6}+\frac9{512g^8},\\
 K_{G_0}=\kappa_{G_0}\sqrt{B(g)},\qquad
 D_{G_0}=\kappa_{G_0}\sqrt{4g^4+216g^2+12}.
 \end{gathered}
 \tag{MR19}
\end{equation}
All quantities are finite positive expressions in the retained
original moments, Gamma weights and $g$. For an endomorphism
$A$ of the cubic coefficient space,
\[
 \|A\|_{G_0}
 =\|G_0^{1/2}AG_0^{-1/2}\|_2
 \le\sqrt{\lambda_{\max}(G_0)\lambda_{\max}(G_0^{-1})}
                 \,\|A\|_{\mathrm F}
 \le\kappa_{G_0}\|A\|_{\mathrm F}.
\]
The first equality follows from the isometric map
$z\mapsto G_0^{1/2}z$. The first inequality follows by
submultiplication and the bound of the operator norm by the
sum of squared matrix entries. The second follows because
each positive eigenvalue is at most the sum of all eigenvalues.
For $\operatorname{diag}(G_0,G_0)$ the largest eigenvalue and
the largest inverse eigenvalue are unchanged, so the same
$\kappa_{G_0}$ applies to the full eight-dimensional operator.
The matrices $G_0$ and $G_0^{-1}$ remain those in RS24--25.

For $0<w\le g$, the absolute value of each entry of $A_w,B_w$
in MR14 is bounded by the corresponding entry below:
\begin{equation}
 A^\flat=
 \begin{pmatrix}
 0&\frac1{16g}&0&\frac3{16}\\
 \frac1{2g^2}&0&\frac1{16g}&0\\
 0&0&0&\frac1{16g}\\
 \frac1{4g^3}&0&0&0
 \end{pmatrix},\qquad
 B^\flat=
 \begin{pmatrix}
 \frac1{4g^3}&0&0&0\\
 0&\frac1{4g^3}&0&0\\
 \frac3{32g^4}&0&0&0\\
 0&\frac3{32g^4}&0&0
 \end{pmatrix}.
 \tag{MR20}
\end{equation}
To verify every bound, use $g+w\ge g$, $2g+w\le3g$ and
$3g+w\le4g$. The entries $\pm1/(16g)$ are unchanged,
the $(0,3)$ entry of $A_w$ is at most $3/16$, its $(1,0)$
entry is at most $1/(2g^2)$, and its $(3,0)$ entry is at most
$1/(4g^3)$. The first two nonzero entries of $B_w$ have
the same last bound, and the remaining two are at most
$3/(32g^4)$. This accounts for every nonzero entry.
Their squared sums are
\[
 \|A^\flat\|_{\mathrm F}^2
 =\frac9{256}+\frac3{256g^2}+\frac1{4g^4}+\frac1{16g^6},
 \qquad
 \|B^\flat\|_{\mathrm F}^2
 =\frac1{8g^6}+\frac9{512g^8}.
\]
Together they equal $B(g)$ in MR19. Since $|c|=1$, MR13
and the proved metric inequality yield the uniform estimate
\begin{equation}
 \left\|\Theta_w^{-1}-\frac{\mathscr L}{w}
       \right\|_{\operatorname{diag}(G_0,G_0)}
 \le K_{G_0},\qquad 0<w\le g.
 \tag{MR21}
\end{equation}
The phase $c$ remains the exact one in MR13; only its
modulus enters this inequality. MR6--7 give the identical
inequality after conjugation by the actual $Q_x$, in
$\mathbb H_U$, with
$\mathscr L_{\rm mark}=Q_x^{-1}\mathscr LQ_x$.

We also need the finite nonzero singular values of the
forward derivative multiplier, rather than only the six
values of the inverse pole. Subtracting the two full
matrices RS2 gives
\begin{equation}
 \frac{H_w-H_0}{w}=
 \begin{pmatrix}
 0&2(2g+w)&0&2(-g^2+gw+w^2)\\
 2&0&2(2g+w)&0\\
 0&-2&0&2(6g-w)\\
 0&0&-2&0
 \end{pmatrix}.
 \tag{MR22}
\end{equation}
This is an identity for every $w>0$, obtained by multiplying
the companion polynomial in RS2 before subtracting its value
at zero. On $0\le w\le g$,
\[
 -g^2\le -g^2+gw+w^2\le g^2.
\]
The lower inequality is $gw+w^2\ge0$; the upper is
$(g-w)(2g+w)\ge0$. Hence the absolute entries of MR22 are
bounded by
\[
 \begin{pmatrix}
 0&6g&0&2g^2\\2&0&6g&0\\0&2&0&12g\\0&0&2&0
 \end{pmatrix}.
\]
Its squared entry sum is exactly $4g^4+216g^2+12$.
Applying the metric estimate from MR19 proves
\begin{equation}
 \|H_w-H_0\|_{G_0}\le wD_{G_0},\qquad
 |\sigma_j^{G_0}(H_w)-\eta_j|\le wD_{G_0},
 \quad j=1,2,\quad 0<w\le g.
 \tag{MR23}
\end{equation}
The second inequality is the same singular-value variational
argument used in MR17; $\eta_1\ge\eta_2>0$ are the two
complete RS8 values in the same $G_0$.

Let $c_6>0$ be the smallest of the six explicitly specified
RS14 constants. Define the following positive threshold:
\begin{equation}
 w_*=\min\left\{
 g,\ \frac{\eta_2}{2D_{G_0}},\
 \frac{c_6}{2(K_{G_0}+1/\eta_2)}
 \right\},\qquad \epsilon_w=wD_{G_0}.
 \tag{MR24}
\end{equation}
This retains the original $w=\delta^2$ and $g=\gamma^2$;
every entry of the threshold is already a finite expression
in the original source data. For $0<w\le w_*$,
$\epsilon_w\le\eta_2/2$. By MR23 the first two singular
values of $H_w$ therefore lie in the positive intervals
$[\eta_j-\epsilon_w,\eta_j+\epsilon_w]$.
Their inverse values $1/(2\sigma_j^{G_0}(H_w))$ are both
at most $1/\eta_2$.

The upper-right block of $\Theta_w^{-1}$ is
$cH_w^{-1}/2$, while its lower-left block is
$cH_w^{-2}/2$. Because the two summands have the same
original Gram $G_0$, their singular spectra combine to
the singular spectrum of the whole inverse, as proved in
RS15. Thus these two bounded reciprocal values occur among
the eight actual singular values. On the other hand MR21
and the variational inequality give
\[
 \sigma_6(\Theta_w^{-1})
 \ge c_6/w-K_{G_0}
 \ge K_{G_0}+2/\eta_2>1/\eta_2.
\]
There are therefore exactly six larger values before those
two reciprocal values. This proves their indices, including
possible equality between $\eta_1$ and $\eta_2$:
\begin{equation}
 \begin{aligned}
 \sigma_7(\Theta_{{\rm mark},w}^{-1})
   &=\frac1{2\sigma_2^{G_0}(H_w)},&
 \left|\sigma_7-\frac1{2\eta_2}\right|
   &\le\frac{wD_{G_0}}{\eta_2^2},\\
 \sigma_8(\Theta_{{\rm mark},w}^{-1})
   &=\frac1{2\sigma_1^{G_0}(H_w)},&
 \left|\sigma_8-\frac1{2\eta_1}\right|
   &\le\frac{wD_{G_0}}{\eta_1^2},\\
 \left|\sigma_j(\Theta_{{\rm mark},w}^{-1})-\frac{c_j}{w}\right|
   &\le K_{G_0},&&1\le j\le6,\qquad 0<w\le w_*.
 \end{aligned}
 \tag{MR25}
\end{equation}
For the finite errors, if $|s-\eta|\le\epsilon_w$ then
\[
 \left|\frac1{2s}-\frac1{2\eta}\right|
 =\frac{|s-\eta|}{2\eta s}
 \le\frac{\epsilon_w}{2\eta(\eta-\epsilon_w)}
 \le\frac{\epsilon_w}{\eta^2}.
\]
The last step uses $\epsilon_w\le\eta_2/2\le\eta/2$.
MR7 transfers both the indexing and inequalities to the
full marked metric, retaining its $E,O$ cross terms.

These estimates also give uniformly positive finite-parameter
exterior bounds. For $0<w\le w_*$ put
\[
 P_r^\pm(w)=\prod_{j=1}^r(c_j\pm wK_{G_0}),
 \qquad 1\le r\le6.
\]
Every minus factor is positive: MR24 gives
$c_j/w-K_{G_0}\ge c_6/w-K_{G_0}>1/\eta_2>0$.
For $1\le r\le6$ and $s=1,2$ the complete bounds are
\begin{equation}
 \begin{aligned}
 \frac{P_r^-(w)}{w^r}
 &\le\|\wedge^r\Theta_{{\rm mark},w}^{-1}\|_{\mathbb H_U}
 \le\frac{P_r^+(w)}{w^r},\\
 \frac{P_6^-(w)}
 {w^6\,2^s\prod_{j=1}^s(\eta_{3-j}+\epsilon_w)}
 &\le\|\wedge^{6+s}\Theta_{{\rm mark},w}^{-1}\|_{\mathbb H_U}
 \le
 \frac{P_6^+(w)}
 {w^6\,2^s\prod_{j=1}^s(\eta_{3-j}-\epsilon_w)} .
 \end{aligned}
 \tag{MR26}
\end{equation}
Every denominator on both sides is positive by MR24.
The first line multiplies the first six inequalities
MR25. The second also uses the exact reciprocal intervals
$[1/(2(\eta_j+\epsilon_w)),1/(2(\eta_j-\epsilon_w))]$
proved above, in their established index order.
Rank zero still has norm one. At rank eight the exact
determinant in MR18 remains the sharper identity
$1/(2^{32}g^6w^6(g+w)^6)$. Thus MR25--26 control both
finite singular directions and all six divergent directions,
with the original source Gram evaluated through every
moment and with its complete marked return.
